\clearpage
\begin{examplebox}
  \textbf{\textcolor{CExample}{例 1：根号当成一个数，二次函数立刻出现（基础）}}

  \textbf{题目：} 求 $y=2x+\sqrt{x-1}$ 的值域。

  \textbf{第一步：找出 $x$ 的范围。} 根号内 $x-1\ge0$，所以 $x\ge1$。

  \textbf{第二步：整个根号设成 $t$。} 令 $t=\sqrt{x-1}$，则 $t\ge0$、$x=t^2+1$，因此
  \[
    y=2(t^2+1)+t=2t^2+t+2.
  \]

  \textbf{第三步：直接比较大小。} 因为 $t\ge0$，所以 $2t^2\ge0$、$t\ge0$，从而
  \[
    y=2+2t^2+t\ge2.
  \]
  当 $t=0$（即 $x=1$）时，恰好取到 $y=2$；随着 $t$ 增大，$y$ 不断增大，没有上限。

  \textbf{答案：}\ $\boxed{[2,+\infty)}$。

  \textbf{特别提醒：} 这里不必先配方。看到 $2t^2+t$ 的两项都不小于 $0$，比配方法更快。

  \medskip
  \textbf{\textcolor{CExample}{例 2：根号里是 $4-x^2$，用圆和三角函数（常规）}}

  \textbf{题目：} 求 $y=x+\sqrt{4-x^2}$ 的值域。

  \textbf{第一步：先定范围。} $4-x^2\ge0$，所以 $-2\le x\le2$。

  \textbf{第二步：用正弦表示 $x$。} 令
  \[
    x=2\sin\theta,\qquad
    -\frac{\pi}{2}\le\theta\le\frac{\pi}{2}.
  \]
  这一段 $\cos\theta\ge0$，所以
  \[
    \sqrt{4-x^2}=2\cos\theta,
    \qquad y=2\sin\theta+2\cos\theta.
  \]

  \textbf{第三步：把两项合成一项。}
  \[
    y=2\sqrt2\sin\left(\theta+\frac{\pi}{4}\right).
  \]
  由于 $-\frac\pi4\le\theta+\frac\pi4\le\frac{3\pi}{4}$，这段正弦的最小值为 $-\frac{\sqrt2}{2}$，最大值为 $1$。所以
  \[
    -2\le y\le2\sqrt2.
  \]
  左端 $-2$ 在 $x=-2$ 时取到；右端 $2\sqrt2$ 在 $x=\sqrt2$ 时取到，中间的值也都能取到。

  \textbf{答案：}\ $\boxed{[-2,\,2\sqrt2]}$。

  \textbf{眼前一亮：} 因为 $\sqrt{4-x^2}\ge0$，所以 $y\ge x\ge-2$，下界无需计算；上界还可用
  \[
    0\le\left(x-\sqrt{4-x^2}\right)^2
      =4-2x\sqrt{4-x^2}
  \]
  得 $x\sqrt{4-x^2}\le2$，继而
  \[
    y^2=x^2+(4-x^2)+2x\sqrt{4-x^2}\le8.
  \]
  因此 $y\le2\sqrt2$，和三角函数方法的答案一致。

  \medskip
  \textbf{\textcolor{CExample}{例 3：原稿中的综合题，换个角度看更简单}}

  \textbf{题目：} 求 $y=\dfrac{\sqrt{x^2+1}}{x-1}$（$x>1$）的值域。

  \textbf{方法 A（使用三角函数）。} 设 $x=\tan\theta$。由于 $x>1$，可取
  \[
    \frac{\pi}{4}<\theta<\frac{\pi}{2}.
  \]
  这段角的余弦为正，于是
  \[
    \begin{aligned}
      y&=\frac{\sqrt{1+\tan^2\theta}}{\tan\theta-1}
         =\frac{1/\cos\theta}{\sin\theta/\cos\theta-1}\\
       &=\frac{1}{\sin\theta-\cos\theta}.
    \end{aligned}
  \]
  当 $\theta$ 从 $\frac\pi4$ 向 $\frac\pi2$ 增大时，$\sin\theta-\cos\theta$ 从接近 $0$（但大于 $0$）增至接近 $1$（但小于 $1$）。因此 $y>1$，且任意大于 $1$ 的值都能取得。

  \textbf{方法 B（更快，不必换三角函数）。} 令 $t=\dfrac{1}{x-1}>0$，则 $x=1+\dfrac1t$。注意 $y>0$，两边平方：
  \[
    \begin{aligned}
      y^2&=\frac{x^2+1}{(x-1)^2}
      =1+\frac{2x}{(x-1)^2}\\
      &=1+2t+2t^2.
    \end{aligned}
  \]
  因为 $t>0$，所以 $y^2>1$、$y>1$。而 $1+2t+2t^2$ 随正数 $t$ 从接近 $0$ 不断增大到没有上限，能够取遍所有大于 $1$ 的数，因此 $y$ 也能取遍所有大于 $1$ 的数。

  \textbf{答案：}\ $\boxed{(1,+\infty)}$。

  \textbf{提醒：} $x$ 永远大于 $1$，不能等于 $1$；即使函数值可以无限接近 $1$，也不能把 $1$ 写入值域。

  \medskip
  \textbf{\textcolor{CExample}{例 4：平方和中，乘一下就能去掉根号（巧妙）}}

  \textbf{题目：} 求 $y=\sqrt{x^2+1}-x$（$x\in\mathbb{R}$）的值域。

  \textbf{第一步：用平方差公式。}
  \[
    y=\frac{(\sqrt{x^2+1}-x)(\sqrt{x^2+1}+x)}{\sqrt{x^2+1}+x}
     =\frac{1}{\sqrt{x^2+1}+x}.
  \]
  因为 $\sqrt{x^2+1}>|x|$，分母始终为正。

  \textbf{第二步：把分母设成 $t$。} 令 $t=\sqrt{x^2+1}+x>0$，于是 $y=\dfrac1t$。

  \textbf{第三步：确认正数 $t$ 全都能取到。} 根据平方差公式，
  \[
    \sqrt{x^2+1}-x=\frac1t,
    \qquad x=\frac{t-1/t}{2}.
  \]
  右式说明每个 $t>0$ 都对应一个实数 $x$，没有遗漏。既然 $t$ 能取遍正数，$\dfrac1t$ 也能取遍正数。

  \textbf{答案：}\ $\boxed{(0,+\infty)}$。

  \textbf{记住：} $\sqrt{x^2+1}-x$ 看起来有减法，结果却始终为正；\textbf{关键就在于根号比 $|x|$ 大。}

  \medskip
  \textbf{\textcolor{CExample}{例 5：平方差根式，反手变成分式（提高）}}

  \textbf{题目：} 求 $y=x-\sqrt{x^2-4}$（$x\ge2$）的值域。

  \textbf{第一步：乘上相加的一半。}
  \[
    y=\frac{\bigl(x-\sqrt{x^2-4}\bigr)\bigl(x+\sqrt{x^2-4}\bigr)}{x+\sqrt{x^2-4}}
     =\frac{4}{x+\sqrt{x^2-4}}.
  \]

  \textbf{第二步：设 $t=x+\sqrt{x^2-4}$。} 因为 $x\ge2$，所以 $t\ge2$，于是 $y=\dfrac4t$。

  \textbf{第三步：检查 $t\ge2$ 是否全部能取到。} 根据两式相加可得
  \[
    x=\frac{t^2+4}{2t}\quad(t\ge2).
  \]
  这个 $x$ 始终不小于 $2$，所以每个 $t\ge2$ 都合法。

  \textbf{第四步：求值域。} 因为 $t\ge2$，故 $0<\dfrac4t\le2$；$t=2$（即 $x=2$）时取到 $2$，而 $0$ 永远取不到。

  \textbf{答案：}\ $\boxed{(0,2]}$。

  \textbf{注意：} 这一题已经给出 $x\ge2$，只需考虑正半边；若原题只写 $\sqrt{x^2-4}$ 的定义域，就还要考虑 $x\le-2$。

  \medskip
  \begin{minipage}{\linewidth}
\textbf{\textcolor{CExample}{例 6：原题限制会缩小换元后的范围（防错）}}
\begin{center}\begin{tikzpicture}[font=\small,>=stealth,line join=round]
\draw[gray!45,rounded corners=3pt] (0,0) rectangle (14.2,9.4);
\node at (7.1,8.8) {换元后的二次函数，只能在允许的新范围上求最值};

\begin{scope}[shift={(4,2.4)},x=1.4cm,y=1cm]
\draw[->] (-.8,0)--(2.9,0) node[right]{$t$};\draw[->] (0,-.3)--(0,5.2) node[above]{$y$};
\draw[gray!50,thick,domain=-.2:2.6,samples=70] plot(\x,{(\x-2)*(\x-2)});\draw[blue,very thick,domain=1:2] plot(\x,{(\x-2)*(\x-2)});
\fill[blue] (1,1) circle(2pt);\fill[blue] (2,0) circle(2pt);\draw[dashed] (1,0)--(1,1)--(0,1);\node[below] at (1,0){$1$};\node[below] at (2,0){$2$};\node[left] at (0,1){$1$};
\end{scope}
\node[align=left,anchor=west] at (8.3,5.1){$2\le x\le5$\\[9pt] $t=\sqrt{x-1}\in[1,2]$\\[9pt] $y=(t-2)^2$\\[9pt] 值域 $[0,1]$};
\node[text width=13.2cm,align=center] at (7.1,.65) {只保留粗蓝段；灰色曲线属于二次式，却不都属于原题。};
\end{tikzpicture}
\end{center}
\end{minipage}



  \textbf{题目：} 求 $y=(\sqrt{x-1}-2)^2$，其中 $2\le x\le5$，的值域。

  令 $t=\sqrt{x-1}$。由于 $2\le x\le5$，有
  \[
    1\le t\le2,\qquad y=(t-2)^2.
  \]
  此时 $-1\le t-2\le0$，所以 $0\le(t-2)^2\le1$。在 $t=2$ 时取 $0$，在 $t=1$ 时取 $1$。

  \textbf{答案：}\ $\boxed{[0,1]}$。

  \textbf{对照：} 如果只写 $t\ge0$，把题目中的 $2\le x\le5$ 忘掉，就会得出完全错误的结果。
\end{examplebox}
